\documentclass[11pt]{article}
\usepackage{fontspec,amsmath,amssymb,amsthm,hyperref}
\usepackage[margin=1in]{geometry}
\setlength{\emergencystretch}{3em}
\title{Sharp Gibbs concentration by dilation}
\author{LangevinLean mathematical exposition}
\date{}
\begin{document}\maketitle
\noindent Source module: \texttt{ProximalLangevin/SharpConcentration.lean}. Compile this proof document with XeLaTeX or LuaLaTeX. The accompanying HTML pairs the mathematical argument with the exact Lean statements and records the scope of the formalization.


Let \(F\) be continuous and \(q\)-strongly convex on a finite-dimensional Hilbert space of dimension \(d\), with \(q>0\) and global minimizer \(z\). Gaussian domination supplies the integrals used below.


\section{A radial contraction compares partition integrals}


Strong convexity and minimality imply \(F(z+r(x-z))\le F(x)-\frac q2(1-r^2)\|x-z\|^2\) for \(0<r<1\). Exponentiate and use \(e^u\ge1+u\). Integrating and changing variables under \(x\mapsto z+r(x-z)\), whose Jacobian is \(r^d\), yields

\[Z+\frac q2(1-r^2)M\le r^{-d}Z,\qquad Z=\int e^{-F},\quad M=\int\|x-z\|^2e^{-F}.\]


Let \(w(x)=e^{-F(x)}\), \(Z=\int w\), \(M=\int\|x-z\|^2w(x)\,dx\), and \(d=\dim H\). Continuity and positive strong convexity first give integrability of both \(w\) and \(\|x-z\|^2w\). Strong convexity and the minimizer growth bound combine to give \[F(z+r(x-z))\le F(x)-\tfrac q2(1-r^2)\|x-z\|^2.\] Exponentiation and \(e^u\ge1+u\) imply \[w(z+r(x-z))\ge w(x)+\tfrac q2(1-r^2)\|x-z\|^2w(x).\]



The local fact \texttt{hpoint} is exactly this last inequality. \texttt{hip} separately proves that the dilated weight is integrable, using translation invariance and nonzero scalar dilation. \texttt{htrans} computes its integral as \(r^{-d}Z\). Integrating \texttt{hpoint} now proves \(Z+\frac q2(1-r^2)M\le r^{-d}Z\), with every integral finite before the comparison.




\textbf{Correspondence with the displayed declarations.}


\begin{itemize}

\item \texttt{sharp\_gibbs\_mode\_dilation\_inequality}: The dilation inequality before taking the tangent limit.
\end{itemize}


\noindent Lean declarations: \texttt{sharp\_gibbs\_mode\_dilation\_inequality}.

\section{Take the sharp limit and normalize}


Multiply the dilation inequality by \(r^d\) and define \(\Phi(r)=r^d[Z+(q/2)(1-r^2)M]\). Then \(\Phi(r)\le Z=\Phi(1)\) for \(0<r<1\). Its left difference quotients at one are nonnegative. Differentiating only this scalar polynomial expression gives \(\Phi'(1)=dZ-qM\ge0\), so \(M\le dZ/q\). The partition function \(Z\) is finite and strictly positive, as proved by strong-convexity Gaussian domination. The actual Gibbs normalization gives

\[\int\|x-z\|^2\pi(dx)=M/Z\le d/q,\qquad\pi(dx)=Z^{-1}e^{-F(x)}dx.\]

The proof differentiates no potential and remains valid for continuous nonsmooth strongly convex \(F\).




Here \texttt{hZ} proves \(Z>0\), and \texttt{hpart} identifies the extended partition function with \(\operatorname{ofReal}Z\). The density change in \texttt{heq} therefore gives \[\int\|x-z\|^2\,d\pi(x)=Z^{-1}M.\] Multiplication of the unnormalized inequality by the nonnegative \(Z^{-1}\), followed by \(Z^{-1}Z=1\), yields \(\int\|x-z\|^2d\pi\le d/q\). Positive normalization is proved from the potential assumptions, rather than supplied as a moment or probability premise.




\textbf{Correspondence with the displayed declarations.}


\begin{itemize}

\item \texttt{sharp\_integral\_sq\_dist\_gibbs\_le}: A continuous \texttt{q}-strongly convex Gibbs potential concentrates around its minimizer with sharp second moment at most \texttt{dim / q}. Every normalization and integrability fact used in the conclusion is proved from the stated inputs.
\end{itemize}


\noindent Lean declarations: \texttt{sharp\_integral\_sq\_dist\_gibbs\_le}.

\section{The one-sided tangent calculation gives d/q}


Retain \(Z,M,d\) from the dilation inequality and define \[\Phi(r)=r^d\bigl[Z+\tfrac q2(1-r^2)M\bigr].\] For \(0<r<1\), dilation proves \(\Phi(r)\le Z=\Phi(1)\). Thus for \(-1<u<0\), the quotient \((\Phi(1+u)-\Phi(1))/u\) is nonnegative: both numerator and denominator are nonpositive.



Differentiate this explicit scalar polynomial, rather than differentiating \(F\): \[\Phi'(1)=dZ-qM.\] In Lean, \texttt{hinner} differentiates the bracket, \texttt{hder} applies the product rule, and \texttt{hnonneg} records the sign of the left difference quotient. Its limit is \(\Phi'(1)\), so \texttt{hbound} gives \(dZ-qM\ge0\). Since \(q>0\), divide to obtain \(M\le(d/q)Z\).




\textbf{Correspondence with the displayed declarations.}


\begin{itemize}

\item \texttt{sharp\_integral\_sq\_dist\_mul\_exp\_neg\_le}: Taking the tangent limit of the dilation inequality gives the exact dimension dependence, without differentiating the potential.
\end{itemize}


\noindent Lean declarations: \texttt{sharp\_integral\_sq\_dist\_mul\_exp\_neg\_le}.

\end{document}
